\documentclass[tikz, border=5pt]{standalone}
\usepackage{xcolor}
\usetikzlibrary{shapes,arrows.meta,calc}

\begin{document}
% Macro para el bloque 3D (toma el color base, p. ej. "gray")
\newcommand{\drawblock}[5]{
	\begin{scope}[shift={(#1)}, rotate=#2]
		\filldraw[fill=#3!30, draw=black] (0,0) rectangle (#4,#5);
		\filldraw[fill=#3!50, draw=black] (#4,0) -- ++(30:0.4) -- ++(0,#5) -- (#4,#5) -- cycle;
		\filldraw[fill=#3!70, draw=black] (0,#5) -- ++(30:0.4) -- ++(#4,0) -- (#4,#5) -- cycle;
	\end{scope}
}
\begin{tikzpicture}[>=Latex, font=\sffamily\fontsize{8pt}{10pt}\selectfont]

	% Coordenadas base
	\coordinate (O) at (0,0);
	\coordinate (laser_start) at (-3,0.3);
	\coordinate (splitter_center) at (-1.5,0.3);
	\coordinate (target_front) at (4.5,0.3);
	\coordinate (receiver_front) at (-1.5,-1.2);

	% Laser Box (3D) -> Pasas "gray" simple, sin signos de admiración
	\coordinate (laser_corner) at ($(laser_start)+(-2.1,-0.3)$);
	\drawblock{laser_corner}{0}{gray}{2.1}{0.6}
	\node[anchor=center, text width=1.5cm, align=center] at ($(laser_start)+(-1.05,0)$) {\bf Láser};

	% Tilted Beam Splitter
	\draw[rotate around={-60:(splitter_center)}, line width=2pt, draw=black!80]
	($(splitter_center)+(0,-0.6)$) -- ($(splitter_center)+(0,0.6)$);

	% Target Block (3D) -> Pasas "gray" simple
	\coordinate (target_corner) at ($(target_front)+(0,-1.2)$);
	\drawblock{target_corner}{0}{gray}{0.8}{2.4}
	\node[anchor=center, align=center] at ($(target_front)+(0.4,-1.8)$) {Superficie};
	\fill[black] (target_front) circle (0.15);

	% Receiver Block (3D) -> Pasas "gray" simple
	\coordinate (receiver_corner) at ($(receiver_front)+(-1.2,-0.8)$);
	\drawblock{receiver_corner}{0}{gray}{2.4}{0.8}
	\node[anchor=center, text width=2cm, align=center] at ($(receiver_front)+(0,-0.4)$) {\bf Receptor};

	% --- Trayectorias de haz ---
	\draw[red, line width=1.2pt, ->] (laser_start) -- (splitter_center);
	\draw[red, line width=1.2pt] (splitter_center) -- (target_front);

	\coordinate (return_start) at ($(target_front)+(0,-0.08)$);
	\coordinate (return_mid) at ($(splitter_center)+(0,-0.08)$);
	\draw[red, line width=1.2pt] (return_start) -- (return_mid);
	\draw[red, line width=1.2pt, ->] (return_start) -- ($(return_mid)!0.5!(return_start)$);

	\draw[red, line width=1.2pt, ->] (return_mid) -- (receiver_front);

	% --- Acotación (distancia R) ---
	\draw[thick] ($(splitter_center)+(0,0.6)$) -- ++(0,1.2);
	\draw[thick] ($(target_front)+(0,1.2)$) -- ++(0,0.6);
	\draw[<->, thick] ($(splitter_center)+(0,1.5)$) -- node[above, font=\bf] {R} ($(target_front)+(0,1.5)$);

\end{tikzpicture}
\end{document}
